\setcounter{section}{18}






  \zwischenueberschrift{Übungsaufgaben}

\inputaufgabe
{}
{





Ein Geldfälscher stellt $3$- und $7$-Euro-Scheine her.
\aufzaehlungdrei{Zeige, dass es nur endlich viele Beträge gibt, die er nicht
\zusatzklammer {exakt} {} {}
begleichen kann. Was ist der höchste Betrag, den er nicht begleichen kann?
}{Was ist der kleinste Betrag, den er auf zwei verschiedene Weisen begleichen kann?
}{Beschreibe die Menge $M$ der vollen Eurobeträge, die er mit seinen Scheinen begleichen kann.
}





}
{} {}

\inputaufgabe
{}
{





Ein Geldfälscher stellt $4$-, $9$- und $11$-Euro-Scheine her.
\aufzaehlungdrei{Zeige, dass es nur endlich viele Beträge gibt, die er nicht
\zusatzklammer {exakt} {} {}
begleichen kann. Was ist der höchste Betrag, den er nicht begleichen kann?
}{Was ist der kleinste Betrag, den er auf zwei verschiedene Weisen begleichen kann?
}{Beschreibe explizit die Menge $M$ der vollen Eurobeträge, die er mit seinen Scheinen begleichen kann.
}





}
{} {}

\inputaufgabegibtloesung
{}
{





Ein Geldfälscher stellt $7$-, $11$-, $13$- und $37$-Euro-Scheine her. Wie viele volle Eurobeträge kann er mit seinen Scheinen nicht bezahlen, und was ist der größte Betrag, den er nicht begleichen kann? Bestimme die
\definitionsverweis {Multiplizität}{}{}
und die
\definitionsverweis {Einbettungsdimension}{}{}
des zugehörigen numerischen Monoids.





}
{} {}

\inputaufgabegibtloesung
{}
{





Bestimme für das
\definitionsverweis {numerische Monoid}{}{}

\mavergleichskette
{\vergleichskette
{ M
}
{ \subseteq }{ \N
}
{  }{
}
{    }{
}
{   }{
}
}
{}{}{,}
das durch $4,7$ und $17$ erzeugt wird, die
\definitionsverweis {Einbettungsdimension}{}{,}
die
\definitionsverweis {Multiplizität}{}{,}
die
\definitionsverweis {Führungszahl}{}{}
und den
\definitionsverweis {Singularitätsgrad}{}{.}





}
{} {}

\inputaufgabe
{}
{





Bestimme für das numerische Monoid $M$, das durch $5,7$ und $9$ erzeugt wird, die
\definitionsverweis {Einbettungsdimension}{}{,} die
\definitionsverweis {Multiplizität}{}{,} die
\definitionsverweis {Führungszahl}{}{} und den
\definitionsverweis {Singularitätsgrad}{}{.}





}
{} {}

\inputaufgabe
{}
{





Es sei

\mavergleichskette
{\vergleichskette
{ M
}
{ \subseteq }{ \N
}
{  }{
}
{    }{
}
{   }{
}
}
{}{}{}
ein
\definitionsverweis {numerisches Monoid}{}{,}
das von teilerfremden natürlichen Zahlen erzeugt sei. Zeige, dass die
\definitionsverweis {Einbettungsdimension}{}{}
maximal gleich der
\definitionsverweis {Multiplizität
}{}{}
ist.





}
{} {}

\inputaufgabe
{}
{





Es sei

\mavergleichskette
{\vergleichskette
{ M
}
{ \subseteq }{ \N
}
{  }{
}
{    }{
}
{   }{
}
}
{}{}{}
ein durch
\definitionsverweis {teilerfremde}{}{}
Zahlen erzeugtes
\definitionsverweis {numerisches Monoid}{}{,}
bei dem die
\definitionsverweis {Einbettungsdimension}{}{}
gleich der
\definitionsverweis {Multiplizität}{}{}
ist. Zeige, dass dann der maximale Erzeuger aus einem minimalen Erzeugendensystem größer oder gleich der
\definitionsverweis {Führungszahl}{}{}
ist.





}
{} {}

\inputaufgabe
{}
{





Man gebe ein Beispiel eines
\definitionsverweis {numerischen Monoids}{}{}
$M$ mit
\definitionsverweis {Multiplizität}{}{}
$3$ und
\definitionsverweis {Einbettungsdimension}{}{}
$3$ an, bei dem die
\definitionsverweis {Führungszahl}{}{}
prim ist und nicht zum minimalen
\definitionsverweis {Erzeugendensystem}{}{}
gehört.





}
{} {}

\inputaufgabe
{}
{





Es sei

\mavergleichskette
{\vergleichskette
{ M
}
{ \subseteq }{ \N
}
{  }{
}
{    }{
}
{   }{
}
}
{}{}{}
ein
\definitionsverweis {numerisches Monoid}{}{,}
das von
\definitionsverweis {teilerfremden}{}{}
Elementen erzeugt werde. Es sei vorausgesetzt, dass die
\definitionsverweis {Multiplizität}{}{}
von $M$ mit der
\definitionsverweis {Führungszahl}{}{}
von $M$ übereinstimmt. Bestimme ein minimales Erzeugendensystem und die
\definitionsverweis {Einbettungsdimension}{}{}
von $M$.





}
{} {}

\inputaufgabegibtloesung
{}
{





Wir betrachten die Neilsche Parabel

\mavergleichskettedisp
{\vergleichskette
{ C
}
{ =} { V \left(  Y^2-X^3  \right)
}
{ \subseteq} { {\mathbb A}^{2}_{ {\mathbb C} }
}
{ } {
}
{ } {
}
}
{}{}{}
und den Punkt

\mavergleichskettedisp
{\vergleichskette
{ P
}
{ =} { (1,1)
}
{ \in} { C
}
{ } {
}
{ } {
}
}
{}{}{.}
Zeige, dass man das maximale Ideal zu $P$, also das Ideal
\mathl{(X-1,Y-1)}{} im Koordinatenring

\mavergleichskette
{\vergleichskette
{ R
}
{ = }{ {\mathbb C} [X,Y]/ \left(  Y^2-X^3  \right)
}
{  }{
}
{    }{
}
{   }{
}
}
{}{}{,}
nicht als
\definitionsverweis {Radikal}{}{}
zu einem einzigen Element

\mavergleichskette
{\vergleichskette
{ f
}
{ \in }{ R
}
{  }{
}
{    }{
}
{   }{
}
}
{}{}{}
beschreiben kann.





}
{} {}

Die vorstehende Aufgabe zeigt, dass jede Kurve $D$, die die Neilsche Parabel im Punkt
\mathl{(1,1)}{} trifft, sie noch in mindestens einem weiteren Punkt trifft. Dies verallgemeinert auch

Aufgabe 1.13
.

\inputaufgabegibtloesung
{}
{





Wir betrachten das von $2$ und $3$
\definitionsverweis {erzeugte Untermonoid}{}{}
in den natürlichen Zahlen, also

\mavergleichskettedisp
{\vergleichskette
{ M
}
{ =} { \{ 0, 2,3,4,5, \ldots \}
}
{ \subset} { \N
}
{ } {
}
{ } {
}
}
{}{}{,}
und den
\definitionsverweis {Restklassenring}{}{}

\mavergleichskette
{\vergleichskette
{ R
}
{ = }{  \Z/( 9)
}
{  }{
}
{    }{
}
{   }{
}
}
{}{}{.}
Ferner betrachten wir die Menge der $R$-wertigen Punkte von
\mathkor {} {M} {und} {\N} {}
\zusatzklammer {also die Menge der Monoidhomomorphismen \mathlk{\operatorname{Mor}_{ \operatorname{mon} } \, ( M ,  R)}{} bzw. \mathlk{\operatorname{Mor}_{ \operatorname{mon} } \, ( \N ,  R)}{}} {} {}
und die Restriktionsabbildung
\maabbeledisp {\psi} { \operatorname{Mor}_{ \operatorname{mon} } \, ( \N ,  R) } {  \operatorname{Mor}_{ \operatorname{mon} } \, ( M ,  R)
} { \pi } {  \pi \vert_M
} {.}
\aufzaehlungsechsabc{Bestimme die
\definitionsverweis {Einheiten}{}{}
und die
\definitionsverweis {nilpotenten}{}{}
Elemente von $R$.
}{Bestimme die Anzahl von
\mathl{\operatorname{Mor}_{ \operatorname{mon} } \, ( \N ,  R)}{.}
}{Zeige, dass ein

\mavergleichskette
{\vergleichskette
{ \rho
}
{ \in }{ \operatorname{Mor}_{ \operatorname{mon} } \, ( M ,  R)
}
{  }{
}
{    }{
}
{   }{
}
}
{}{}{}
mit der Eigenschaft, dass $\rho(2)$ eine Einheit in $R$ ist, eine Fortsetzung nach $\N$ besitzt.
}{Bestimme sämtliche

\mavergleichskette
{\vergleichskette
{ \pi
}
{ \in }{ \operatorname{Mor}_{ \operatorname{mon} } \, ( \N ,  R)
}
{  }{
}
{    }{
}
{   }{
}
}
{}{}{,}
deren Einschränkung auf $M \setminus \{0\}$ die Nullfunktion ist.
}{Bestimme sämtliche

\mavergleichskette
{\vergleichskette
{ \rho
}
{ \in }{ \operatorname{Mor}_{ \operatorname{mon} } \, ( M ,  R)
}
{  }{
}
{    }{
}
{   }{
}
}
{}{}{,}
die keine Fortsetzung nach $\N$ besitzen.
}{Bestimme die Anzahl von
\mathl{\operatorname{Mor}_{ \operatorname{mon} } \, ( M ,  R)}{.}
}





}
{} {}

\inputaufgabe
{}
{





Es sei $M$ ein numerisches Monoid. Bestimme die
\definitionsverweis {Filter}{}{}
in $M$.





}
{} {}

\inputaufgabe
{}
{





Es sei $M$ ein
\definitionsverweis {numerisches Monoid}{}{}
und $K$ ein
\definitionsverweis {Körper}{}{.}
Zeige, dass es ein

\mavergleichskette
{\vergleichskette
{ f
}
{ \in }{ K[M]
}
{  }{
}
{    }{
}
{   }{
}
}
{}{}{}
derart gibt, dass

\mavergleichskette
{\vergleichskette
{ K[M]_f
}
{ = }{ K[X]_X
}
{  }{
}
{    }{
}
{   }{
}
}
{}{}{}
gilt.





}
{} {}

\inputaufgabe
{}
{





Es sei $M$ ein
\definitionsverweis {numerisches Monoid}{}{,}
$K$ ein
\definitionsverweis {Körper}{}{}
und
\mathl{K[M]}{} der zugehörige
\definitionsverweis {Monoidring}{}{.}
Zeige, dass

\mavergleichskette
{\vergleichskette
{ K[M]
}
{ \cong }{ K[\N]
}
{  }{
}
{    }{
}
{   }{
}
}
{}{}{}
genau dann gilt, wenn

\mavergleichskette
{\vergleichskette
{M
}
{ = }{\N
}
{  }{
}
{    }{
}
{   }{
}
}
{}{}{}
ist.





}
{} {}

\inputaufgabegibtloesung
{}
{





Man gebe ein Beispiel für einen
\definitionsverweis {injektiven}{}{}
\definitionsverweis {Monoidhomomorphismus}{}{}
\maabb {} { M } { N
} {}
zwischen kommutativen
\definitionsverweis {Monoiden}{}{}
derart an, dass die zugehörige
\definitionsverweis {Spektrumsabbildung}{}{}
nicht
\definitionsverweis {surjektiv}{}{}
ist.





}
{} {}

\inputaufgabe
{}
{





Es sei $M$ ein
\definitionsverweis {kommutatives Monoid}{}{}
und sei $K$ ein
\definitionsverweis {Körper}{}{.}
\aufzaehlungdrei{Zeige, dass durch die
\definitionsverweis {Diagonalabbildung}{}{}
\maabbeledisp {\triangle} { M } { M \times M
} { m } { (m,m)
} {,}
ein
\definitionsverweis {Monoidhomomorphismus}{}{}
gegeben ist.
}{Beschreibe den zugehörigen
$K$-\definitionsverweis {Algebrahomomorphismus}{}{}
\maabbdisp {} {K[M] } { K[M \times M] \cong K[M] \otimes K[M]
} {.}
}{Beschreibe die zugehörige
\definitionsverweis {Spektrumsabbildung}{}{}
\maabbdisp {} { K\!-\!\operatorname{Spek}\, \left(  K[M]   \right) \times K\!-\!\operatorname{Spek}\, \left(  K[M]   \right)  } { K\!-\!\operatorname{Spek}\, \left(  K[M]   \right)
} {.}
Zeige insbesondere, dass dadurch
\mathl{K\!-\!\operatorname{Spek}\, \left(  K[M]   \right)}{} selbst zu einem kommutativen Monoid wird.
}





}
{} {}

\inputaufgabe
{}
{





Spezialisiere

Aufgabe 18.16

auf die Monoide

\mavergleichskettedisp
{\vergleichskette
{M
}
{ =} { \N, \Z, \N^r, \Z^r
}
{ } {
}
{ } {
}
{ } {
}
}
{}{}{.}





}
{} {}

\inputaufgabe
{}
{





Es sei $M$ ein
\definitionsverweis {kommutatives Monoid}{}{,}
$K$ ein
\definitionsverweis {Körper}{}{}
und

\mavergleichskette
{\vergleichskette
{ P
}
{ \in }{ K\!-\!\operatorname{Spek}\, \left(   M   \right)
}
{  }{
}
{    }{
}
{   }{
}
}
{}{}{}
ein fixierter Punkt.
\aufzaehlungdreiabc{Zeige, dass es eine
\definitionsverweis {stetige Abbildung}{}{}
\maabbeledisp {\Psi_P} { K\!-\!\operatorname{Spek}\, \left(   M   \right) } { K\!-\!\operatorname{Spek}\, \left(   M   \right)
} { Q} { Q \cdot P
} {,}
gibt.
}{Beschreibe den zugehörigen
$K$-\definitionsverweis {Algebrahomomorphismus}{}{}
\maabb {} { K[M] } { K[M]
} {.}
}{Charakterisiere, für welche $P$ diese Abbildung
\definitionsverweis {bijektiv}{}{}
ist.
}





}
{} {}

\inputaufgabe
{}
{





Es seien $M$ und $N$
\definitionsverweis {kommutative Monoide}{}{,}
sei
\maabb {\varphi} { M } { N
} {}
ein
\definitionsverweis {Monoidhomomorphismus}{}{}
und sei $K$ ein
\definitionsverweis {Körper}{}{.}
Zeige, dass die
\definitionsverweis {Spektrumsabbildung}{}{}
\maabbdisp {\varphi^*} { K\!-\!\operatorname{Spek}\, \left(   N   \right) } { K\!-\!\operatorname{Spek}\, \left(   M   \right)
} {}
ebenfalls ein Monoidhomomorphismus ist.





}
{} {}

\inputaufgabe
{}
{





Es seien

\mavergleichskette
{\vergleichskette
{ a,b,c,d,r,s
}
{ \geq }{ 1
}
{  }{
}
{    }{
}
{   }{
}
}
{}{}{}
\definitionsverweis {natürliche Zahlen}{}{.}
Wir betrachten die
\definitionsverweis {stetig differenzierbare Kurve}{}{}
\maabbeledisp {} { [0,1] } { \R^2
} { t } {   \left(  t^r   , \,   t^s                   \right)
} {.}
Berechne das
\definitionsverweis {Wegintegral}{}{}
längs dieses Weges zum
\definitionsverweis {Vektorfeld}{}{}

\mavergleichskettedisp
{\vergleichskette
{ F(x,y)
}
{ =} {  \left(  x^ay^b   , \,   x^c y^d                   \right)
}
{ } {
}
{ } {
}
{ } {
}
}
{}{}{.}





}
{} {}





  \zwischenueberschrift{Aufgaben zum Abgeben}

\inputaufgabe
{6}
{





Es sei $M$ ein
\definitionsverweis {numerisches Monoid}{}{,}
das durch zwei
\definitionsverweis {teilerfremde}{}{}
Elemente

\mavergleichskette
{\vergleichskette
{ d
}
{ > }{ e
}
{  }{
}
{    }{
}
{   }{
}
}
{}{}{}
erzeugt werde. Bestimme die
\definitionsverweis {Einbettungsdimension}{}{,}
die
\definitionsverweis {Multiplizität}{}{,}
die
\definitionsverweis {Führungszahl}{}{}
und den
\definitionsverweis {Singularitätsgrad}{}{}
von $M$.





}
{} {}

\inputaufgabe
{3}
{





Bestimme für das numerische Monoid $M$, das durch
\mathl{3,7,9}{} und $11$ erzeugt wird, die
\definitionsverweis {Einbettungsdimension}{}{,} die
\definitionsverweis {Multiplizität}{}{,} die
\definitionsverweis {Führungszahl}{}{} und den
\definitionsverweis {Singularitätsgrad}{}{.}





}
{} {}

\inputaufgabe
{4}
{





Klassifiziere sämtliche
\definitionsverweis {numerische Monoide}{}{}
$M$
\zusatzklammer {mit teilerfremden Erzeugern} {} {} mit
\definitionsverweis {Führungszahl}{}{}

\mavergleichskette
{\vergleichskette
{ f(M)
}
{ \leq }{ 6
}
{  }{
}
{    }{
}
{   }{
}
}
{}{}{.}
Man gebe jeweils die
\definitionsverweis {Einbettungsdimension}{}{,}
die
\definitionsverweis {Multiplizität}{}{}
und den
\definitionsverweis {Singularitätsgrad}{}{}
an.





}
{} {}

\inputaufgabe
{8 (1+3+1+2+1)}
{





Wir betrachten das von $2$ und $5$
\definitionsverweis {erzeugte Untermonoid}{}{}
in den natürlichen Zahlen, also

\mavergleichskettedisp
{\vergleichskette
{ M
}
{ =} { \{ 0, 2,4,5,6, \ldots \}
}
{ \subset} { \N
}
{ } {
}
{ } {
}
}
{}{}{,}
und den
\definitionsverweis {Restklassenring}{}{}

\mavergleichskette
{\vergleichskette
{ R
}
{ = }{  \Z/( 4)
}
{  }{
}
{    }{
}
{   }{
}
}
{}{}{.}
Ferner betrachten wir die Menge der $R$-wertigen Punkte von
\mathkor {} {M} {und} {\N} {}
\zusatzklammer {also die Menge der Monoidhomomorphismen \mathlk{\operatorname{Mor}_{ \operatorname{mon} } \, ( M ,  R)}{} bzw. \mathlk{\operatorname{Mor}_{ \operatorname{mon} } \, ( \N ,  R)}{}} {} {}
und die Restriktionsabbildung
\maabbeledisp {\psi} { \operatorname{Mor}_{ \operatorname{mon} } \, ( \N ,  R) } {  \operatorname{Mor}_{ \operatorname{mon} } \, ( M ,  R)
} { \pi } {  \pi \vert_M
} {.}
\aufzaehlungfuenfabc{Bestimme die
\definitionsverweis {Einheiten}{}{}
und die
\definitionsverweis {nilpotenten}{}{}
Elemente von $R$ und die Anzahl von
\mathl{\operatorname{Mor}_{ \operatorname{mon} } \, ( \N ,  R)}{.}
}{Zeige, dass ein

\mavergleichskette
{\vergleichskette
{ \rho
}
{ \in }{ \operatorname{Mor}_{ \operatorname{mon} } \, ( M ,  R)
}
{  }{
}
{    }{
}
{   }{
}
}
{}{}{}
mit der Eigenschaft, dass $\rho(2)$ eine Einheit in $R$ ist, eine Fortsetzung nach $\N$ besitzt.
}{Bestimme sämtliche

\mavergleichskette
{\vergleichskette
{ \pi
}
{ \in }{ \operatorname{Mor}_{ \operatorname{mon} } \, ( \N ,  R)
}
{  }{
}
{    }{
}
{   }{
}
}
{}{}{,}
deren Einschränkung auf $M \setminus \{0\}$ die Nullfunktion ist.
}{Bestimme sämtliche

\mavergleichskette
{\vergleichskette
{ \rho
}
{ \in }{ \operatorname{Mor}_{ \operatorname{mon} } \, ( M ,  R)
}
{  }{
}
{    }{
}
{   }{
}
}
{}{}{,}
die keine Fortsetzung nach $\N$ besitzen.
}{Bestimme die Anzahl von
\mathl{\operatorname{Mor}_{ \operatorname{mon} } \, ( M ,  R)}{.}
}





}
{} {}

\inputaufgabe
{3}
{





Es sei

\mavergleichskette
{\vergleichskette
{ M
}
{ \subseteq }{ \N
}
{  }{
}
{    }{
}
{   }{
}
}
{}{}{}
ein
\definitionsverweis {numerisches Monoid}{}{}
und $K$ ein
\definitionsverweis {Körper}{}{.}
Definiere

\mathdisp {M_+= M \cap \N_+} {  }

und

\mathdisp {n M_+ = \left\{ m \in M  \mid  \text{es gibt eine Darstellung } m = m_1 + \cdots + m_n  \text{ mit } m_i \in M_+        \right\}} { . }

Zeige, dass $nM_+$ \anfuehrung{Ideale}{} in $M$ sind, dass zu $M_+$ ein
\definitionsverweis {maximales Ideal}{}{}
${\mathfrak m}$ in $K[M]$ gehört, und dass das zu $nM_+$ gehörige Ideal gleich ${\mathfrak m}^n$ ist.





}
{} {}

\inputaufgabe
{4}
{





Es sei $K$ ein
\definitionsverweis {algebraisch abgeschlossener Körper}{}{}
und seien $M,N$
\definitionsverweis {numerische Monoide}{}{}
mit

\mavergleichskette
{\vergleichskette
{ M
}
{ \subseteq }{ N
}
{  }{
}
{    }{
}
{   }{
}
}
{}{}{.}
Zeige, dass die zugehörige
\definitionsverweis {Spektrumsabbildung}{}{}
\definitionsverweis {surjektiv}{}{}
ist.





}
{Es ist dabei hilfreich,

Satz 18.10

zu verwenden.} {}

\inputaufgabe
{3}
{





Es seien $M,N$
\definitionsverweis {numerische Monoide}{}{.}
Für welche der numerischen Invarianten $\nu$
\zusatzklammer {\definitionsverweis {Multiplizität}{}{,}
\definitionsverweis {Führungszahl}{}{,}
\definitionsverweis {Singularitätsgrad}{}{,}
\definitionsverweis {Einbettungsdimension}{}{}} {} {}
folgt aus

\mavergleichskette
{\vergleichskette
{ M
}
{ \subseteq }{ N
}
{  }{
}
{    }{
}
{   }{
}
}
{}{}{}
die Abschätzung

\mavergleichskette
{\vergleichskette
{ \nu(M)
}
{ \geq }{ \nu(N)
}
{  }{
}
{    }{
}
{   }{
}
}
{}{}{?}





}
{(Beweis oder Gegenbeispiel!)} {}

\inputaufgabe
{3}
{





Es sei $M$ ein
\definitionsverweis {numerisches Monoid}{}{,}
das nicht isomorph zu $\N$ sei, und sei $K$ ein Körper. Zeige, dass es im Monoidring $K[M]$ \definitionsverweis {irreduzible Elemente}{}{}
gibt, die nicht
\definitionsverweis {prim}{}{}
sind. Man gebe Elemente aus $K[M]$ mit zwei wesentlich verschiedenen Zerlegungen in irreduzible Elemente an.





}
{} {}
